% answers_short.tex —— 1977 年·黑龙江初试卷 · 简短答案速查
% 由 tools/make_short_answers.py 自 answers.tex 的 \gksolhead/\gkanswer 提取，
% 并与 1977/1977_数学真题_黑龙江初试_答案要点.md 逐题比对后生成，勿手改（改源 MD 后重跑）。

\begin{gknotice}
\textbf{使用说明}：本节只给最终答案便于快速核对，完整推导见解析版。答案为编写者独立推导，\textbf{非官方标准答案}。
\end{gknotice}

\gkpart{题目}
\begin{itemize}[leftmargin=2.6em]
  \item[\textbf{1.}] $\dfrac{13}{15}$。
  \item[\textbf{2.}] $6\dfrac59$。
  \item[\textbf{3.}] 半径 $r=6\ \mathrm{cm}$，圆的面积 $S=113.04\ \mathrm{cm^2}$。
  \item[\textbf{4.}] $x=-5$。
  \item[\textbf{5.}] $3(a+b-2c)(a+b+2c)$。
  \item[\textbf{6.}] 见解析版
  \item[\textbf{7.}] $\odot O$ 的半径为 $5\sqrt3\ \mathrm{cm}$，切线 $AB=5\ \mathrm{cm}$。
  \item[\textbf{8.}] 见解析版
  \item[\textbf{9.}] 甲车间原计划生产 $200$ 件产品，乙车间原计划生产 $250$ 件产品。
  \item[\textbf{10.}] ① $-\dfrac12$；② $-1$。
  \item[\textbf{11.}] $BC=3\sqrt6\ \mathrm{cm}$。
  \item[\textbf{12.}] 最大值为 $4$（当 $x=-1$ 时取得）。
\end{itemize}
